% =====================================================================
%  IEMS5726 Data Science in Practice  (2025-26 Term 2 / Session A)
%  Group 43 — Project Report
%  Compile with: pdflatex on Overleaf (article class, A4)
% =====================================================================
\documentclass[11pt,a4paper]{article}

% ----- packages ------------------------------------------------------
\usepackage[utf8]{inputenc}
\usepackage[T1]{fontenc}
\usepackage{lmodern}
\usepackage{amsmath}
\usepackage[margin=2.4cm]{geometry}
\usepackage{graphicx}
\usepackage{xcolor}
\usepackage{hyperref}
\usepackage{titlesec}
\usepackage{enumitem}
\usepackage{caption}
\usepackage{float}
\usepackage{placeins}
\usepackage{microtype}
\usepackage{parskip}
\usepackage{booktabs}

% ----- colors --------------------------------------------------------
\definecolor{sdNavy}{HTML}{1A1A4A}
\definecolor{sdGrey}{HTML}{555555}

% ----- section heading style ----------------------------------------
\titleformat{\section}{\Large\bfseries\color{sdNavy}}{\thesection.}{0.6em}{}
\titleformat{\subsection}{\large\bfseries\color{sdNavy}}{\thesubsection.}{0.5em}{}
\titlespacing*{\section}{0pt}{1.0em}{0.55em}
\titlespacing*{\subsection}{0pt}{0.75em}{0.35em}

% ----- start every new section on a new page -------------------------
% FloatBarrier prevents figures from drifting into the next section.
\newcommand{\sectionbreak}{\FloatBarrier\clearpage}

% ----- caption style -------------------------------------------------
\captionsetup{
  font={small,it,color=sdGrey},
  labelfont={small,it,color=sdGrey},
  justification=centering,
  labelsep=period,
  skip=4pt
}

% ----- paragraph and list style -------------------------------------
\setlength{\parskip}{5pt}
\setlength{\parindent}{0pt}

\setlist[itemize]{leftmargin=1.35em,topsep=2pt,itemsep=2pt,parsep=0pt}
\setlist[enumerate]{leftmargin=1.55em,topsep=2pt,itemsep=3pt,parsep=0pt}

% ----- hyperref defaults --------------------------------------------
\hypersetup{
  pdftitle={SonicDNA -- A Multi-Lens Music Personality Mirror},
  pdfauthor={Group 43 (Xupeng ZHANG, Zetao HUANG)},
  colorlinks=true,
  linkcolor=sdNavy,
  urlcolor=sdNavy,
  citecolor=sdNavy
}

% ----- shorthand: normal figure -------------------------------------
% The image is constrained by BOTH width and height.
% This prevents wide screenshots from extending beyond the right margin.
\newcommand{\sdfigure}[3][0.88\linewidth]{%
  \begin{figure}[!htbp]
    \centering
    \includegraphics[width=#1,height=0.60\textheight,keepaspectratio]{#2}
    \caption{#3}
  \end{figure}
}

% ----- shorthand: large screenshot figure ----------------------------
% Use this for UI screenshots. It still cannot exceed text width.
\newcommand{\sdscreenshot}[3][0.92\linewidth]{%
  \begin{figure}[!htbp]
    \centering
    \includegraphics[width=#1,height=0.62\textheight,keepaspectratio]{#2}
    \caption{#3}
  \end{figure}
}

% ----- shorthand: small/narrow figure --------------------------------
\newcommand{\sdsmallfigure}[3][0.48\linewidth]{%
  \begin{figure}[!htbp]
    \centering
    \includegraphics[width=#1,height=0.55\textheight,keepaspectratio]{#2}
    \caption{#3}
  \end{figure}
}

% =====================================================================
\begin{document}

% =====================================================================
% Title page
% =====================================================================
\begin{titlepage}
  \thispagestyle{empty}
  \centering

  % ----- top course information -----
  \vspace*{1.6cm}

  {\Large\bfseries IEMS5726 Data Science in Practice}\\[6pt]
  {\normalsize 2025--26 Term 2 / Session A}

  % ----- controlled vertical gap -----
  \vspace{2.6cm}

  % ----- main title block -----
  {\fontsize{26pt}{31pt}\selectfont\bfseries
   SonicDNA}\\[10pt]

  {\fontsize{18pt}{22pt}\selectfont\bfseries
   A Multi-Lens Music Personality Mirror}\\[1.4cm]

  {\large\bfseries Project Report}

  % ----- separator -----
  \vspace{1.2cm}
  \rule{0.68\linewidth}{0.45pt}
  \vspace{1.0cm}

  % ----- group information -----
  \begin{tabular}{rl}
    \textbf{Group number:} & 43 \\[5pt]
    \textbf{Group members:} & 1155238738, Xupeng ZHANG \\[3pt]
                            & 1155251352, Zetao HUANG \\
  \end{tabular}

  % ----- bottom date -----
  \vfill

  {\normalsize\textit{Submission date: 8 May 2026}}

  \vspace*{1.2cm}

\end{titlepage}

% =====================================================================
\tableofcontents

% =====================================================================
\section{Problem Definitions}

Modern music streaming services typically expose only two surfaces to a
listener: an algorithmic ``similar tracks'' rail and an annual
marketing recap (e.g. Spotify Wrapped). Both reduce the user's library
to a single dimension --- popularity or genre --- and leave no room for
self-exploration of taste.

In short, SonicDNA turns a playlist into a visual, psychological and
recommendation-based profile. SonicDNA addresses the gap with a
\emph{multi-lens music personality mirror}: the user uploads a playlist
(CSV) and the application produces three complementary outputs derived
from a single 9-dimensional audio-feature embedding:

\begin{itemize}
  \item \textbf{MoodPrint} --- a vinyl-record artwork (a ``music DNA'')
        deterministically derived from the playlist composition, so two
        different playlists never share a print.
  \item \textbf{Music MBTI} --- the playlist is reduced to a 4-letter
        code on four interpretable axes (E/I, N/S, F/T, J/P), each axis
        built from a subset of audio features, mapping to one of 16
        named personality types with hand-written 60-character
        descriptions.
  \item \textbf{Parallel-Universe Playlists} --- for three hand-curated
        ``universes'' (Neon Synthwave Drive / Lo-fi Study Den /
        Mosh Pit Riot) we linearly project the user's centroid toward
        the universe's centroid ($\alpha = 0.7$) and retrieve the top-10
        most cosine-similar tracks inside that universe, answering
        ``what would you listen to if you grew up there?''.
\end{itemize}

The application therefore solves a dual problem: it produces objective,
reproducible analytics (clustering, PCA, recommendation) and a
subjective interpretive layer (MoodPrint art + MBTI labels) in one
coherent pipeline --- making the analytical results memorable and
shareable instead of disposable.

% =====================================================================
\section{Data Collection, Preprocessing and Representation}

\subsection{Data source}

Kaggle \emph{Spotify Tracks Dataset} (Maharshi Pandya).

\begin{sloppypar}
\noindent\url{https://www.kaggle.com/datasets/maharshipandya/-spotify-tracks-dataset}
\end{sloppypar}

The dataset contains 114{,}000 tracks across 114 fine-grained genres,
with the full audio-features panel that the Spotify Web API used to
expose. Using this offline dump lets us avoid the API access changes
published 2024-11 by Spotify.

\subsection{Data collection}

We fetch the dataset programmatically via \texttt{kagglehub} on first
run (see \path{preprocess.py} under \path{source_code/pipeline/} and
the \texttt{Dockerfile} build step), so the application is
self-contained: no external API, no long-running scrapers.

\subsection{Data preprocessing}

\begin{itemize}
  \item Drop the original index column \texttt{Unnamed: 0}.
  \item Drop rows missing any of \{\texttt{track\_id}, \texttt{artists},
        \texttt{track\_name}\} (3 rows).
  \item De-duplicate on \texttt{track\_id} (the dataset duplicates each
        track once per genre slice; 114k rows $\to$ 89{,}740 unique
        tracks).
  \item Aggregate the 114 fine genres into 11 meta-genres (Pop, Rock,
        Metal, HipHop, Electronic, Jazz, Folk, Latin, World, Chill,
        Niche) via a hand-built dictionary in
        \path{pipeline/preprocess.py}.
  \item Standardize the 9 numeric audio features with sklearn's
        \texttt{StandardScaler} (mean 0, std 1) so cosine similarity,
        KMeans and PCA all operate in a comparable feature space.
  \item Persist the cleaned table to parquet (snappy) and the scaler to
        joblib pickle for re-use by all downstream modules.
\end{itemize}

\subsection{Data representation}

Each track is represented as a 9-dimensional real-valued vector in
z-scored space:
\begin{center}
  \emph{(danceability, energy, valence, acousticness, instrumentalness,
   liveness, speechiness, tempo, popularity).}
\end{center}

A user playlist is represented as the mean of these 9-vectors across
its tracks. This single embedding is the input shared by MoodPrint,
Music MBTI, the recommender, and the parallel-universe projection ---
keeping all four creative modules consistent and letting any
improvement to the representation propagate downstream without code
changes elsewhere.

% =====================================================================
\section{Data Modeling}

\subsection{Mood KMeans (k = 8)}

We fit one KMeans (\texttt{k=8, random\_state=42, n\_init=10}) on the
z-scored 9-D feature space. Cluster names are derived from each
centroid's most distinctive trait via a small rule-base (see
\path{pipeline/cluster.py::name_mood_clusters}): Mainstream Pulse,
Sunlit Indie, Library Hush, Live Stage Glow, Stormcore Drive, Cipher
Booth, Rainy Window, Drift Layer.

\sdfigure[0.82\linewidth]{figures/05_centroid_radars.png}{%
  KMeans 8 cluster centroids in z-scored space. Each cluster has a
  distinct dominant axis (e.g.\ Cipher Booth on speechiness; Library
  Hush on acousticness and instrumentalness).}

\subsection{Music MBTI (4-axis composite mapping)}

On the same z-scored space we define four composite axes whose signs
encode an MBTI-style 4-letter code:

\begin{itemize}
  \item \textbf{E vs I}: energy + danceability (extrovert vs.\
        introvert).
  \item \textbf{N vs S}: acousticness + instrumentalness (intuitive
        vs.\ sensing).
  \item \textbf{F vs T}: $\text{valence} - 0.5 \times
        \text{speechiness}$ (feeling vs.\ thinking).
  \item \textbf{J vs P}: popularity (mainstream-judging vs.\
        niche-perceiving).
\end{itemize}

Per-track axes are also computed
(\path{pipeline/mbti.py::attach_track_axes}), which lets us list
``representative tracks'' for each of the 16 types by picking the
tracks whose own 4-axis scores are deepest into their quadrant. This
avoids the sample-imbalance trap of training a 16-class classifier on
heavily skewed labels.

\subsection{Parallel Universe projection}

For each of three hand-curated universes we filter the library to a
curated set of fine-grained genres (e.g., Neon Synthwave Drive =
\{synth-pop, electronic, electro, disco, new-age, house, deep-house,
progressive-house, chicago-house, edm\}) and compute the centroid
vector. The user's centroid $\mathbf{u}$ is then linearly interpolated
toward each universe centroid $\mathbf{a}$:
\[
  \mathbf{p} \;=\; \mathbf{u} + \alpha (\mathbf{a} - \mathbf{u}),
  \qquad \alpha = 0.7.
\]
We run cosine-similarity recommendation inside that universe's track
pool only, so the result is a believable ``you, but in this universe''
playlist instead of a hard genre swap that would lose all of the
user's personal style.

\subsection{Diversity-aware recommender}

Pure cosine ranking suffers from the all-tracks-by-the-same-artist
failure mode. We post-process the cosine ordering with two hard
diversity caps: at most 1 track per primary artist and at most 3
tracks per meta-genre. This is the same routine reused by both the
global recommender (Dashboard tab) and the universe-projected
recommender (Parallel Universes tab) via a
\texttt{library\_subset\_mask}.

\subsection{Pretrained model artifacts}

All sklearn artifacts are saved to \path{source_code/models/}:

\begin{itemize}
  \item \texttt{scaler.pkl} --- StandardScaler over 9 audio features
  \item \texttt{kmeans\_mood.pkl} --- KMeans($k=8$) over z-scored space
  \item \texttt{pca.pkl} --- 2-D PCA used by the Insights scatter map
  \item \texttt{mood\_cluster\_names.json} --- interpretive name per
        centroid id
\end{itemize}

\begin{sloppypar}
\noindent The model artifacts are included in the accompanying
submission package, and the source repository at
\url{https://github.com/ShepZhang/sonicdna-iems5726} also includes
a \path{pipeline/preprocess.py} script that regenerates them from
the Kaggle dataset in roughly two minutes.
\end{sloppypar}

% =====================================================================
\section{Data Visualization}

\subsection{Visualization for the data}

\sdfigure[0.82\linewidth]{figures/01_feature_distributions.png}{%
  Distributions of 9 audio features on the raw scale across
  89{,}740 tracks; acousticness / instrumentalness are heavily
  right-skewed, tempo is roughly bimodal around 90 and 130 BPM.}

\sdfigure[0.62\linewidth]{figures/02_correlation_heatmap.png}{%
  Pearson correlations between the 9 modeling features. The set is
  mostly uncorrelated ($|r|<0.5$ everywhere except the energy--loudness
  pair which we do not include), meaning the 9-D vector is not
  redundant.}

\sdfigure[0.80\linewidth]{figures/03_metagenre_distribution.png}{%
  Track count per meta-genre after collapsing 114 fine genres into 11
  buckets.}

\subsection{Visualization for the model}

\sdfigure[0.76\linewidth]{figures/04_pca_global_map.png}{%
  PCA 2-D projection of the full 89{,}740-track library, colored by
  KMeans mood cluster id. The 8 clusters separate cleanly even with
  only 40.5\% of variance preserved.}

\subsection{Visualization for the results}

Each user playlist produces a unique MoodPrint. Below is the MoodPrint
of a 25-track Chill playlist used for development:

\sdsmallfigure[0.50\linewidth]{figures/moodprint_smoke_chill.png}{%
  MoodPrint for a 25-track Chill playlist. The vinyl record encodes 9
  audio features as concentric tracks; the outer ring shows the
  mood-cluster proportions.}

% =====================================================================
\section{System Architecture}

\subsection{Current MVP architecture}

The application is a single Streamlit process backed by a decoupled
pipeline package. All four creative modules live in
\path{source_code/pipeline/} and are imported by
\path{application/app.py} via an explicit \path{sys.path} injection at
the top of the module:

\sdfigure[0.84\linewidth]{figures/06_arch_mvp.png}{%
  Current MVP architecture. Streamlit is the only UI surface; all model
  artifacts live on disk and are loaded once via
  \texttt{@st.cache\_resource}.}

\subsection{Future production architecture}

The pipeline package is intentionally framework-agnostic so the
Streamlit UI can be swapped for the standard production stack
(NGINX $\to$ Vue3 SPA + FastAPI backend $\to$ MySQL/Redis) without
touching any modeling code. The diagram below documents that upgrade
path:

\sdfigure[0.84\linewidth]{figures/07_arch_production.png}{%
  Future production architecture upgrade path. The Streamlit app
  collapses into a Vue3 SPA + FastAPI backend; the same pipeline
  package is reused as a library, so no modeling code changes.}

\subsection{Deployment}

The MVP is fully containerized: a single \texttt{Dockerfile} (build
context = project root) installs both \path{source_code} and
\path{application} dependencies, downloads the Kaggle dataset at build
time, runs the preprocessing + clustering pipeline once, and exposes
the Streamlit app on port 8501. The same image is the deliverable for
the public HuggingFace Spaces demo.

\begin{itemize}
  \item Source repository:
        \url{https://github.com/ShepZhang/sonicdna-iems5726}
  \item Public live demo: provided in the accompanying submission
        package once the HuggingFace Space has been published.
\end{itemize}

% =====================================================================
\section{User Interfaces of the Software}

The Streamlit application has a fixed sidebar (sample-or-upload
playlist selector) and four content tabs.

\paragraph{Screenshot 1 --- Dashboard tab.}
Vinyl-style MoodPrint, Music MBTI card with rarity stat, the 16-type
distribution bar, and the four-axis breakdown.

\sdscreenshot[0.90\linewidth]{screenshots/01_dashboard.png}{%
  Dashboard tab --- Workout Beast Mode sample. The vinyl print encodes
  9 audio features as concentric tracks; arc length, brightness and
  thickness reflect how strongly the playlist exceeds the global
  average on each feature.}

\paragraph{Screenshot 2 --- Parallel Universes tab.}
Segmented control to pick one of three hand-curated universes, Top-10
cosine recommendations from inside that universe's pool, and a
three-overlay radar.

\sdscreenshot[0.90\linewidth]{screenshots/02_universes.png}{%
  Parallel Universes tab with Neon Synthwave Drive selected. The
  yellow polygon is the projected blend ($\alpha = 0.7$) used to score
  recommendations.}

\paragraph{Screenshot 3 --- Insights tab.}
PCA 2-D projection of the full 89{,}740-track library coloured by mood
cluster, with the user's tracks highlighted as pink-ringed dots; below
it are the user's 9-axis taste radar and mood-cluster donut.

\sdscreenshot[0.90\linewidth]{screenshots/03_insights.png}{%
  Insights tab --- PCA scatter with the Workout playlist clustered in
  the high-energy / high-popularity region of the global music map.}

\paragraph{Screenshot 4 --- Sidebar.}
Language-stable English-only controls, sample-vs-upload selector, CSV
schema help tooltip, and a What's-Inside cheatsheet for the four tabs.

\sdsmallfigure[0.42\linewidth]{screenshots/06_sidebar.png}{%
  Sidebar --- sample selector, CSV-upload switch, and a brief tab
  cheatsheet.}

% =====================================================================
\section{Challenges}

Up to four challenges, in order of effort spent:

\begin{enumerate}
  \item \textbf{9-D audio features $\to$ 2-D artistic fingerprint
        (MoodPrint).}\par
        The main difficulty was balancing uniqueness and
        interpretability. Reducing 9 dimensions into a single image
        while preserving both \emph{uniqueness} and \emph{readability}
        required two design choices: (a) fixing each feature's hue so
        the 9 tracks are mnemonically learnable, and (b) standardizing
        the z-score with
        $\mathrm{strength}=\mathrm{clip}((z+2)/4,0,1)$, then mapping it
        to arc length, line width and alpha. As a result, mean
        differences as small as $0.3\sigma$ can still produce visibly
        different rings. The outer 8-cluster ring, sized by the
        playlist's cluster mass, further strengthens the uniqueness of
        the overall print.

  \item \textbf{Sample imbalance in 16-class Music MBTI.}\par
        The main difficulty was avoiding a misleading supervised
        classifier. Naively training a 16-class classifier on the
        library would have produced wildly unbalanced classes (ESFJ has
        12{,}672 tracks; ISFP only 2{,}035). We sidestepped this by
        defining the four MBTI axes as composite z-score sums rather
        than learned weights, making classification deterministic and
        interpretable. Representative tracks are then generated by
        ranking each type's quadrant by total signed axis strength and
        de-duplicating on artist.

  \item \textbf{Plausible ``parallel-universe'' projection without
        losing user identity.}\par
        The main difficulty was changing the musical world while still
        preserving the user's personal taste. A naive implementation
        would replace the user's vector with the universe centroid,
        producing identical recommendations for every user. Instead we
        linearly interpolate
        $\mathbf{u} + \alpha (\mathbf{a} - \mathbf{u})$ with
        $\alpha = 0.7$, preserving 30\% of the user's personal
        direction inside each universe pool. Cosine recommendation is
        then applied inside the pool with the same diversity caps as
        the main recommender.

  \item \textbf{Working around the Spotify Web API 2024-11 access
        policy change.}\par
        The main difficulty was keeping the project reproducible
        without depending on unstable API access. Spotify deprecated
        audio-features endpoints for new development apps. Instead of
        registering a Spotify dev app and depending on external token
        freshness at request time, the entire pipeline operates on the
        offline Kaggle dump fetched once via \texttt{kagglehub} at
        Docker build time. This keeps cold-start time short and makes
        the project resilient to future API changes.
\end{enumerate}

% =====================================================================
\section{References}

\begin{itemize}
  \item Maharshi Pandya. \emph{Spotify Tracks Dataset}, 2023.\\
        \url{https://www.kaggle.com/datasets/maharshipandya/-spotify-tracks-dataset}
  \item Pedregosa et al. scikit-learn: Machine Learning in Python.
        \emph{JMLR}, 12, 2825--2830, 2011.
  \item Hunter, J.~D. Matplotlib: A 2D graphics environment.
        \emph{Computing in Science \& Engineering}, 9(3), 90--95, 2007.
  \item McKinney, W. pandas: a Foundational Python Library for Data
        Analysis and Statistics. \emph{Python for High Performance and
        Scientific Computing}, 14, 1--9, 2011.
  \item Streamlit Inc. \emph{Streamlit --- the fastest way to build
        data apps}. \url{https://streamlit.io}
  \item Plotly Technologies Inc. \emph{Plotly Python Open Source
        Graphing Library}. \url{https://plotly.com/python/}
  \item HuggingFace. \emph{Spaces Documentation}.
        \url{https://huggingface.co/docs/hub/spaces}
  \item Spotify. \emph{Introducing some changes to our Web API}, 27
        November 2024.
        \url{https://developer.spotify.com/blog/2024-11-27-changes-to-the-web-api}
\end{itemize}

% =====================================================================
\section{Declaration for the use of AI Tools}

AI assistants were used in the following capacities during this
project:

\begin{itemize}
  \item Cursor was used to draft and iteratively refine the pipeline
        modules (\path{preprocess.py}, \path{cluster.py},
        \path{moodprint.py}, \path{mbti.py},
        \path{parallel_universe.py}, \path{recommender.py}) and the
        Streamlit application (\path{application/app.py}). All
        algorithmic decisions, including the cluster naming rule book,
        MoodPrint polar layout, MBTI 4-axis definitions, $\alpha=0.7$
        interpolation and diversity caps, were specified by the author;
        the assistant was used to translate those decisions into
        runnable code.
  \item AI assistance was used to generate the 16 short Music-MBTI
        personality descriptions. Each description was reviewed and
        edited by the author before inclusion.
  \item AI was \emph{not} used to select or download the dataset, train
        the models, invent any result numbers, or replace the author's
        final judgement. AI assistance was used to polish the report
        structure and wording.
\end{itemize}

% =====================================================================
\section{Signed VeriGuide Receipt}

The signed VeriGuide receipt is attached as the last page of this
report.

\begin{center}
  \fbox{\parbox{0.82\linewidth}{\centering
    \textcolor{red}{\textbf{[ INSERT VERIGUIDE RECEIPT IMAGE HERE ]}}\\[6pt]
    \small After submitting to VeriGuide and saving the receipt PNG
    as \texttt{veriguide.png} alongside this \texttt{.tex}, replace
    this box with:\\[4pt]
    \texttt{\textbackslash includegraphics[}\\
    \texttt{\quad width=0.92\textbackslash linewidth,}\\
    \texttt{\quad height=0.85\textbackslash textheight,}\\
    \texttt{\quad keepaspectratio]\{veriguide.png\}}}}
\end{center}

% When the final VeriGuide receipt image is ready, replace the box above
% with the following code:
%
% \begin{center}
%   \includegraphics[width=0.92\linewidth,height=0.85\textheight,keepaspectratio]{veriguide.png}
% \end{center}

\end{document}
